% Part: lambda-calculus
% Chapter: lambda-definability
% Section: pairs

\documentclass[../../../include/open-logic-section]{subfiles}

\begin{document}

\olfileid[id]{lam}{ldf}{pai}
\olsection{Pasangan dan Pendahulu}

\begin{defn}
Pasangan $M$ dan $N$ (ditulis $\tuple{M,N}$) didefinisikan sebagai berikut:
\[
\tuple{M,N} \ident \lambd[f][fMN].
\]
\end{defn}
  
Secara intuitif, pasangan merupakan fungsi yang menerima suatu fungsi, lalu
menerapkan fungsi itu pada kedua elemen pasangan. Mengikuti gagasan ini, kita
memperoleh konstruktor berikut, yang menerima dua suku dan mengembalikan
pasangan yang memuat keduanya:
\[
  \fn{Pair} \ident \lambd[mn][\lambd[f][fmn]]
\]
Jika diberikan suatu pasangan, kita juga ingin memperoleh kembali
elemen-elemennya. Untuk itu, kita memerlukan dua fungsi pengakses yang menerima
suatu pasangan sebagai argumen dan mengembalikan elemen pertama atau kedua dari
pasangan itu:
\begin{align*}
  \fn{Fst} & \ident \lambd[p][p(\lambd[mn][m])]\\
  \fn{Snd} & \ident \lambd[p][p(\lambd[mn][n])]
\end{align*}

\begin{prob}
  Jelaskan mengapa fungsi pengakses $\fn{Fst}$ dan $\fn{Snd}$ bekerja.
\end{prob}

Sekarang, dengan pasangan, kita dapat !!{lambda define} fungsi pendahulu:
\[
  \fn{Pred} \ident \lambd[n][\fn{Fst}(n (\lambd[p][\tuple{\fn{Snd}\, {p}, \fn{Succ}(\fn{Snd}\, {p})}]) \tuple{\num 0, \num 0})]
\]
Ingat bahwa $\num n\, f x$ tereduksi menjadi $f^{n}(x)$; dalam hal ini,
$f$ adalah fungsi yang menerima suatu pasangan $p$ dan mengembalikan pasangan
baru yang memuat komponen kedua $p$ dan penerus komponen kedua itu; sedangkan
$x$ adalah pasangan $\tuple{0,0}$. Dengan demikian, hasilnya adalah
$\tuple{0,0}$ jika $n=0$, dan $\tuple{\num{n-1}, \num n}$ dalam kasus
lainnya. Kemudian $\fn{Pred}$ mengembalikan komponen pertama hasil tersebut.

Pengurangan dapat didefinisikan sebagai $\fn{Pred}$ yang diterapkan pada $a$
sebanyak $b$ kali:
\[
\fn{Sub} \ident \lambd[ab][b \fn{Pred}\, a].
\]
    
\end{document}
